\chapter{Background}\label{ch:background}

For the purpose of this thesis, it was necessary for us to go through some of the content from the course COMP9334, which I have compiled below.

\section{Capacity Planning Concepts}

Capacity planning concerns predicting how a system will perform under
different workloads and determining the resources needed to meet service
objectives. Key terminology includes:

\begin{itemize}
    \item \textbf{Throughput} ($X$): the rate at which a system completes work.
    \item \textbf{Service demand} ($D$): the total service time a job
          requires from a resource.
    \item \textbf{Utilisation} ($U$): the fraction of time a resource is busy.
    \item \textbf{Bottleneck device}: the resource with the largest
          service demand, which limits maximum throughput.
    \item \textbf{Saturation}: the point at which utilisation approaches
          100\%, causing rapid growth in response time.
\end{itemize}

These concepts provide the foundation for evaluating system behaviour
before, during, and after load increases.

\section{Queueing Theory Foundations}

Queueing theory models how jobs compete for shared resources. A queue is
typically denoted using \emph{Kendall notation}, such as M/M/1, where:

\begin{itemize}
    \item the first letter describes the arrival process (M = memoryless / Poisson),
    \item the second describes the service-time distribution 
          (M = exponential, G = general),
    \item the final value indicates the number of servers (1, $c$, etc.).
\end{itemize}

For an M/M/1 queue, arrival rate $\lambda$ and service rate $\mu$ yield:

\[
\rho = \frac{\lambda}{\mu}, \qquad R = \frac{1}{\mu - \lambda},
\]

where $\rho$ is utilization and $R$ is mean response time.  
Systems remain stable only when $\rho < 1$.

The M/M/$c$ and M/G/1 models generalize these results to multiple servers
and general service-time distributions. For M/G/1, the Pollaczek–Khinchine
formula characterizes mean waiting time based on the first two moments of
service time.

\section{Operational Laws}

Operational analysis provides model-free relationships that hold for any
stable system:

\begin{itemize}
    \item \textbf{Utilisation Law}:\quad $U = X S$
    \item \textbf{Little’s Law}:\quad $N = X R$
    \item \textbf{Forced Flow Law}:\quad $X(j) = V(j) X$
\end{itemize}

Here $S$ is the mean service time, $R$ is the mean response time, $N$ is the mean
number of jobs in the system, and $V(j)$ is the visit ratio of device $j$.
These laws enable service-demand extraction and consistency checking
using measured performance data. By collecting performance data and creating queueing network models, operational analysis helps identify bottlenecks and forecast system behavior before deployment, improving design decisions.

\section{Discrete-Event Simulation}

Many systems cannot be solved analytically, such as queues with general
inter-arrival distributions, complex scheduling policies, or multiple
interacting resources. In such cases, \textbf{discrete-event simulation
(DES)} provides an alternative.

A DES model tracks the system over time by processing events such as job
arrivals, service completions, and state transitions. Core components of
a DES include:

\begin{itemize}
    \item a simulation clock,
    \item event list (future events),
    \item queues and servers,
    \item random variate generation for arrival and service times.
\end{itemize}

DES is used when:

\begin{itemize}
    \item analytical models (e.g.\ G/G/1) have no closed-form solutions,
    \item service-time or arrival distributions deviate from ideal assumptions,
    \item the system contains multiple interacting bottlenecks or
          resource constraints.
\end{itemize}

While flexible, DES requires a careful choice of parameters, warm-up
periods, run lengths, and variance-reduction techniques to ensure
statistically reliable results.

\section{Previous Work}

Teaching-oriented literature on capacity planning commonly introduces
queuing models, operational laws, and simulation methods as the core
toolkit for analyzing system performance. Foundational texts emphasize
reproducible modeling techniques, service-demand extraction, and
comparison of analytical predictions with simulation results. These
methods form the theoretical basis upon which later practical and
experimental work can be built. Now we are going to examine some of the experimental work and industrial practices that have directed my thesis.
